% Research experience entries — \input after \section{...}
% Uses \roleentry{title}{org}{dates}{advisor}{bullets}
% Use \cvonly{} and \resumeonly{} for content visible in only one doc

% Project sub-heading inside a role: \project{title}{sub-bullets}
\newcommand{\project}[2]{%
  \item \textbf{#1}%
  \begin{itemize}[leftmargin=12pt,topsep=0pt,itemsep=1pt,parsep=0pt]%
  #2%
  \end{itemize}%
}

\roleentry%
  {Undergraduate Researcher}%
  {Neural Systems Laboratory, UW}%
  {Oct 2025--Present}%
  {Advisor: Dr.\ Rajesh Rao, Paul G. Allen School of Computer Science \& Engineering}%
  {%
    \project{msTBFM: multi-session stimulation response modeling (second author, TNSRE)}{%
      \item Co-designed and ran every experiment in the paper: Monte Carlo cross-validation, cross-animal transfer, latency benchmarks, and ablations.
      \item Key results: compared with single-session models, meta-learned pretraining reduced catastrophic forecast failures from 16/40 to 1/40 sessions and calibration data needs by 50--90\%.
      \item Built the testing harness, analysis pipeline, statistical tests, and cloud infrastructure used to run every experiment at scale; designed most results figures.
      \cvonly{\item Built co-adaptation, a non-meta-learning control for training and test-time adaptation, showing MAML's inner--outer loop improves sample efficiency and robustness.}
    }
    \project{Honors thesis: traveling-wave signals in motor cortex, toward closed-loop control}{%
      \item Testing whether wave dynamics give a leading indicator of behavior for closed-loop stimulation; building models to forecast latent cortical dynamics under stimulation.
      \item Decoding lever trajectory, trial outcome, and reaction time from generalized-phase gradients in rodent M1/M2 surface-grid recordings (Purdue collaborators).
      \item Preliminary: pre-cue distributed cortical activity predicts upcoming hits vs.\ missed cues, surviving motor-state controls and spike-only replication.
    }
    \cvonly{\project{Predictive model of the hippocampal--entorhinal system}{%
      \item Extending Active Predictive Coding and the Tolman-Eichenbaum Machine into a two-stream model of the entorhinal cortex and hippocampus, with hypernetwork-generated transitions and cross-stream memory.
      \item Testing for emergent grid-cell-like codes and for planning and recall in BabyAI and a custom Recall Rooms task.
    }}
  }

\roleentry%
  {Associate Researcher, VP Operations}%
  {Sensors, Energy, and Automation Laboratory (SEAL), UW}%
  {Jan 2025--Oct 2025}%
  {Advisor: Dr.\ Alexander Mamishev, Department of Electrical \& Computer Engineering}%
  {%
    \item Developed signal processing for VLF/ULF electromagnetic measurements for geophysical monitoring.
    \item Wrote sections of and coordinated federal and industry grant proposals (DOE STTR, NIH, NSF, CDC, Google).
    \cvonly{\item Led a project team building lab automation and DevOps tooling.}
  }
